\section{Operational Amplifier (Op-Amp)}\label{sec:devices_opamp}

An op-amp is treated here as a device: a three-terminal (plus supplies),
active, (ideally) linear, time-invariant differential voltage amplifier.
How it is built is the subject of Section~\ref{sec:amp_differential}; here
we only need its terminal behaviour to use it inside blocks
(Section~\ref{sec:amp_functional}).

\begin{circuitfig}
  \draw (0,0) node[op amp](A){};
  \draw (A.-) node[left]{$V_-$} (A.+) node[left]{$V_+$} (A.out) node[right]{$V_{out}$};
  \draw (4,1) node[left]{$V_+$} to[short, o-] (4.8,1) to[european resistor, l_={$R_{id}$}] (4.8,-1) to[short, -o] (4,-1) node[left]{$V_-$};
  \draw (9,1) to[european resistor, l={$R_o$}] (10.8,1) to[short, -o] (11.3,1) node[right]{$V_{out}$};
  \draw (9,-1.2) node[ground]{} to[american controlled voltage source, l={$A(s)(V_+ - V_-)$}] (9,1);
\end{circuitfig}

\begin{itemize}
  \item \textbf{Mathematical Model}:
    \[
      V_{out} = A(s) \left( V_+ - V_- \right) = A(s) V_d
    \]
    \begin{itemize}
      \item \textbf{Ideal Model}: \(A \to \infty\), \(R_{id} \to \infty\),
        \(R_o \to 0\). For a finite output with infinite gain,
        \begin{align*}
          V_d &= \frac{V_{out}}{A} \xrightarrow{A \to \infty} 0
          \hspace{0.5cm} \Rightarrow \hspace{0.5cm} V_+ = V_- \hspace{0.5cm} \text{[virtual short, negative feedback only]} \\
          I_+ &= I_- = \frac{V_d}{R_{id}} = 0 \hspace{0.5cm} \text{[no input current]}
        \end{align*}
      \item \textbf{Single Pole Model}: a dominant pole \(p_1\) and DC gain
        \(A_0\),
        \begin{align*}
          A(s) &= \frac{A_0}{1 + \frac{s}{p_1}} = \frac{A_0 p_1}{s + p_1} \\
          GB &= A_0 \cdot \frac{p_1}{2\pi} \hspace{1cm} \text{[gain-bandwidth product, Hz]} \\
          A(s) &\approx \frac{2\pi GB}{s} \hspace{1cm} [|s| \gg p_1]
        \end{align*}
      \item \textbf{Large Signal Limits (Q5/Q6)}: the output saturates and
        slews,
        \begin{align*}
          V_{OL} \leq V_{out} \leq V_{OH} \hspace{1cm}
          \left|\frac{dV_{out}}{dt}\right| \leq SR
        \end{align*}
        A sine of amplitude \(V_p\) and frequency \(f\) has maximum slope
        \(2\pi f V_p\), so it is undistorted only while
        \[
          f \leq f_{FP} = \frac{SR}{2\pi V_p} \hspace{1cm} \text{[full power bandwidth]}
        \]
    \end{itemize}
  \item \textbf{Negative Feedback}: A network returns a fraction \(F\) of
    the output to the inverting input, \(V_- = F V_{out}\), \(V_+ = V_{in}\).
    \begin{align*}
      V_{out} &= A(V_{in} - F V_{out}) \\
      V_{out}(1 + AF) &= A V_{in} \\
      A_{CL} &= \frac{V_{out}}{V_{in}} = \frac{A}{1 + AF} = \frac{1}{F} \cdot \frac{AF}{1 + AF} \xrightarrow{AF \to \infty} \frac{1}{F}
    \end{align*}
    \(T = AF\) is the loop gain. With the single pole model:
    \begin{align*}
      A_{CL}(s) &= \frac{\frac{A_0 p_1}{s + p_1}}{1 + \frac{A_0 p_1 F}{s + p_1}} \\
      A_{CL}(s) &= \frac{A_0 p_1}{s + p_1 + A_0 p_1 F} \\
      A_{CL}(s) &= \frac{A_0 p_1}{s + p_1(1 + A_0 F)} \\
      A_{CL}(s) &= \frac{\frac{A_0}{1 + A_0 F}}{1 + \frac{s}{p_1(1 + A_0 F)}}
    \end{align*}
    The closed loop DC gain falls by \(1 + A_0F\) and the pole rises by
    \(1 + A_0F\): gain times bandwidth is conserved,
    \[
      f_{-3dB,CL} = \frac{p_1 (1 + A_0 F)}{2\pi} \approx F \cdot GB
    \]
  \item \textbf{Non-Idealities}:
    \begin{itemize}
      \item \textbf{Offset} \(V_{OS}\): the \(V_d\) needed to bring
        \(V_{out}\) to its intended value. It appears in series with \(V_+\),
        so it is amplified by \(\frac{1}{F}\).
      \item \textbf{Bias Currents} \(I_{B+}\), \(I_{B-}\) flow into the
        inputs; \(I_{OS} = |I_{B+} - I_{B-}|\).
      \item \textbf{Common Mode Rejection}: with
        \(V_{ic} = \frac{V_+ + V_-}{2}\),
        \begin{align*}
          V_{out} &= A_d V_d + A_c V_{ic} \hspace{1cm}
          CMRR = \frac{A_d}{A_c}
        \end{align*}
      \item \textbf{Power Supply Rejection}:
        \(PSRR = \frac{A_d}{V_{out}/V_{dd}}\), the gain from the input divided
        by the gain from the supply.
    \end{itemize}
  \item \textbf{Circuit Analysis}: In Q5 check that \(V_{out}\) is inside
    \(V_{OL} \ldots V_{OH}\) and the input common mode is inside its range.
    In Q4 write KCL with \(V_{out} = A(s)(V_+ - V_-)\) and take
    \(A \to \infty\) at the end; this keeps the derivation honest about
    which results need high gain.
\end{itemize}
